\documentclass[12pt,a4paper]{article}
\usepackage{amsmath,amssymb,amsthm}
\usepackage{hyperref}
\usepackage{geometry}
\usepackage{enumitem}
\geometry{margin=2.5cm}
\usepackage{booktabs}

\newtheorem{theorem}{Theorem}[section]
\newtheorem{lemma}[theorem]{Lemma}
\newtheorem{corollary}[theorem]{Corollary}
\newtheorem{proposition}[theorem]{Proposition}
\theoremstyle{definition}
\newtheorem{definition}[theorem]{Definition}
\newtheorem*{hypothesis}{Hypothesis}
\theoremstyle{remark}
\newtheorem{remark}[theorem]{Remark}

\title{Nuclear Binding Energy from Recognition Science:\\
Emergent Bethe--Weizs\"acker Structure\\
from the $J$-Cost Functional}

\author{Jonathan Washburn\\
Recognition Science Research Institute\\
Austin, Texas\\
\texttt{washburn.jonathan@gmail.com}}

\date{March 2026}

\begin{document}
\maketitle

\begin{abstract}
We show that the functional form of the Bethe--Weizs\"acker
(BW) mass formula emerges naturally from the Recognition
Science (RS) $J$-cost framework.  The five BW terms---volume,
surface, Coulomb, asymmetry, and pairing---correspond to
distinct ledger-structure aspects: bulk recognition filling,
boundary recognition cost, the derived $\alpha$, the $J$-cost
penalty for neutron-proton imbalance, and $\varphi$-resonant
nucleon pairs.  The qualitative features of nuclear binding
(iron peak at $A \approx 56$, magic numbers, stability line)
follow from the ledger topology.  The asymmetry term is the
most direct RS contribution: the BW asymmetry formula is the
leading-order Taylor expansion of $J(N/Z)$, derived here
without any free parameters.  The precise numerical BW
coefficients remain an open first-principles problem.
\end{abstract}

%% ============================================================
\section{Introduction}
\label{sec:intro}
%% ============================================================

Nuclear binding energy $B(A,Z)$---the energy required to
disassemble a nucleus of mass number $A$ and atomic number $Z$
into its constituent nucleons---is one of the most consequential
quantities in physics.  It determines stellar nucleosynthesis,
energy production (fission and fusion), the iron peak in
elemental abundances, and the stability of all matter.

The semi-empirical Bethe--Weizs\"acker (BW)
formula~\cite{BW1935,Bethe1936} captures the gross features of
$B(A,Z)$ with five terms, each with an empirically fitted
coefficient.  In the Standard Model, these coefficients emerge
from QCD in principle, but a first-principles lattice QCD
calculation of nuclear binding for $A > 4$ remains beyond current
computational reach.

In Recognition Science~\cite{RS_Axioms}, the $J$-cost functional
$J(x) = \tfrac{1}{2}(x + x^{-1}) - 1$ governs recognition
dynamics on the 3D discrete ledger.  We show that the
\emph{functional form} of the BW formula---specifically, the
$A$, $A^{2/3}$, $Z^2 A^{-1/3}$, $(A-2Z)^2/A$, and pairing
dependences---arises naturally from the ledger structure, and
that the energy scale is set by the $\varphi$-ladder.

%% ============================================================
\section{The Bethe--Weizs\"acker Formula}
\label{sec:bw}
%% ============================================================

\begin{theorem}[Nuclear Binding Energy]
\label{thm:binding}
The binding energy of a nucleus $(A, Z)$ is:
\begin{equation}
  \boxed{B(A,Z) = a_V A - a_S A^{2/3} -
  a_C \frac{Z(Z-1)}{A^{1/3}} -
  a_A \frac{(A - 2Z)^2}{A} + \delta(A,Z),}
  \label{eq:bw}
\end{equation}
where the five terms are: volume ($a_V$), surface ($a_S$),
Coulomb ($a_C$), asymmetry ($a_A$), and pairing ($\delta$).
\end{theorem}

The empirical values~\cite{Krane} are:

\begin{center}
\renewcommand{\arraystretch}{1.3}
\begin{tabular}{@{}lcc@{}}
\toprule
\textbf{Coefficient} & \textbf{Value (MeV)} &
\textbf{RS interpretation} \\
\midrule
$a_V$ & $15.75$ & Recognition filling energy \\
$a_S$ & $17.8$ & Boundary recognition cost \\
$a_C$ & $0.711$ & Electromagnetic repulsion
($\propto \alpha$) \\
$a_A$ & $23.7$ & Isospin imbalance cost \\
$\delta_0$ & $11.2/\sqrt{A}$ & $\varphi$-resonant pairing \\
\bottomrule
\end{tabular}
\end{center}

%% ============================================================
\section{RS Origin of Each Term}
\label{sec:terms}
%% ============================================================

\subsection{Volume Term: $a_V A$}

\begin{proposition}[Volume from Ledger Filling]
\label{prop:volume}
When $A$ nucleons occupy a compact region of the ledger, each
nucleon participates in recognition events with its neighbors.
For a nucleon in the bulk (fully surrounded), the recognition
energy gain per nucleon is approximately constant.  The total
binding from recognition filling therefore scales as $A$.
\end{proposition}

\begin{proof}
On the 3D cubic ledger, each voxel has 6 nearest neighbors.
A nucleon at voxel $\mathbf{n}$ recognizes its neighbors via
the $J$-cost: the cost of recognition is $J(x)$ where $x$ is
the ratio of voxel states.  When neighbors have compatible
states (bound nucleons), $x \approx 1$ and $J(x) \approx 0$,
releasing the maximum binding per pair.  For a bulk nucleon,
this gives a binding contribution proportional to the
coordination number, which is independent of $A$ for
$A \gg 1$.  Summing over all nucleons yields $B_V \propto A$.
\end{proof}

\begin{remark}
The energy scale $a_V \approx 15.75$ MeV lies on the
$\varphi$-ladder at $a_V \approx E_{\mathrm{coh}} \cdot
\varphi^{r_V}$ for some nuclear rung $r_V$.  With
$E_{\mathrm{coh}} = \varphi^{-5} \approx 0.0902$ eV,
we need $\varphi^{r_V} \approx 1.75 \times 10^8$,
giving $r_V \approx 39$.  Deriving $r_V$ from first
principles (i.e., from the QCD binding of three quarks into
a proton and the residual strong force between protons)
remains an open problem.
\end{remark}

\subsection{Surface Term: $-a_S A^{2/3}$}

\begin{proposition}[Surface from Boundary Cost]
\label{prop:surface}
Nucleons at the nuclear surface have fewer neighbors than
bulk nucleons, incurring an excess $J$-cost.  The number of
surface nucleons scales as $A^{2/3}$.
\end{proposition}

\begin{proof}
Model the nucleus as a sphere of radius $R \propto A^{1/3}$.
Its surface area $4\pi R^2 \propto A^{2/3}$.  Each surface
nucleon has $\sim 3$ neighbors instead of 6, losing
approximately half its recognition binding.  The total
surface penalty is $\Delta B_S \propto A^{2/3}$.
\end{proof}

\subsection{Coulomb Term: $-a_C Z(Z-1)/A^{1/3}$}

\begin{proposition}[Coulomb from Fine-Structure Constant]
\label{prop:coulomb}
The Coulomb repulsion between $Z$ protons in a sphere of
radius $R \propto A^{1/3}$ yields:
\begin{equation}
  E_C = \frac{3}{5} \frac{Z(Z-1)\alpha}{R}
  \propto \frac{Z(Z-1)}{A^{1/3}}.
\end{equation}
In RS, $\alpha$ is derived from $D = 3$ cube geometry:
$\alpha^{-1} = 4\pi \cdot 11 - w_8 \ln\varphi +
\delta_\kappa \approx 137.036$.
\end{proposition}

\begin{remark}
The Coulomb coefficient $a_C \approx 0.711$ MeV is the only
BW coefficient that RS can compute from first principles,
because $\alpha$ is derived.  The nuclear radius
$R = r_0 A^{1/3}$ with $r_0 \approx 1.2$ fm enters through
the single SI calibration seam.
\end{remark}

\subsection{Asymmetry Term: $-a_A (A-2Z)^2/A$}

\begin{proposition}[$J$-Cost Imbalance]
\label{prop:asymmetry}
The $J$-cost functional penalizes deviations from $x = 1$.
In the nuclear context, the ``ratio'' is the neutron-to-proton
balance $x = N/Z$.  The penalty is:
\begin{equation}
  J(N/Z) = \tfrac{1}{2}\!\left(\frac{N}{Z} + \frac{Z}{N}\right)
  - 1 \geq 0,
\end{equation}
with equality when $N = Z$.  For small asymmetry
$\epsilon = (N-Z)/(N+Z)$:
\begin{equation}
  J(N/Z) \approx 2\epsilon^2 = 2\left(\frac{A - 2Z}{A}\right)^2.
  \label{eq:j-asym}
\end{equation}
Summing over all nucleons:
$E_A \propto J(N/Z) \cdot A \propto (A-2Z)^2/A$.
\end{proposition}

\begin{proof}
Write $N/Z = (1+\epsilon)/(1-\epsilon)$ for
$\epsilon = (N-Z)/(N+Z) = (A-2Z)/(2A)$.  Then:
\begin{align}
  J(N/Z) &= \tfrac{1}{2}\!\left(\frac{1+\epsilon}{1-\epsilon}
  + \frac{1-\epsilon}{1+\epsilon}\right) - 1 \nonumber \\
  &= \frac{(1+\epsilon)^2 + (1-\epsilon)^2}{2(1-\epsilon^2)} - 1
  = \frac{1+\epsilon^2}{1-\epsilon^2} - 1
  = \frac{2\epsilon^2}{1-\epsilon^2}
  \approx 2\epsilon^2 \quad (\epsilon \ll 1).
\end{align}
Multiplying by $A$ (number of nucleon pairs scaling) and
absorbing the factor of 2 into the coefficient $a_A$:
$E_A = a_A \cdot (A-2Z)^2 / A$.
\end{proof}

\begin{remark}
This is one of the most direct connections between the $J$-cost
functional and nuclear physics: the asymmetry penalty is
\emph{literally} the $J$-cost evaluated at the neutron-proton
ratio.  The functional form $(A-2Z)^2/A$ is not imposed by
hand but follows from the Taylor expansion of $J(N/Z)$.
\end{remark}

\subsection{Pairing Term: $\delta(A,Z)$}

\begin{proposition}[Pairing from $\varphi$-Resonance]
\label{prop:pairing}
Nucleon pairs with compatible $\varphi$-ladder states gain
extra binding.  This favors even-even nuclei
($\delta > 0$), penalizes odd-odd nuclei ($\delta < 0$),
and vanishes for odd-$A$ nuclei ($\delta = 0$).
\end{proposition}

The pairing interaction arises because two nucleons in the
same orbital with opposite spin form a $J$-cost minimum at
$x = 1$ (recognition event with zero cost).  Even-even nuclei
have all nucleons paired; odd-odd nuclei have two unpaired
nucleons.  Empirically,
$\delta \approx 11.2 / \sqrt{A}$ MeV.

%% ============================================================
\section{The Iron Peak}
\label{sec:iron}
%% ============================================================

\begin{theorem}[Iron Peak]
\label{thm:iron}
The binding energy per nucleon $B/A$ peaks at
$A \approx 56$--$62$ (the iron--nickel group).
\end{theorem}

\begin{proof}
Setting $\partial(B/A)/\partial A = 0$ for symmetric nuclei
($Z = A/2$) and neglecting pairing:
\begin{equation}
  \frac{\partial}{\partial A}\!\left(a_V - a_S A^{-1/3} -
  \frac{a_C}{4} A^{2/3}\right) = 0,
\end{equation}
which gives:
\begin{equation}
  A_{\mathrm{peak}} = \frac{2\,a_S}{a_C}
  \approx \frac{2 \times 17.8}{0.711}
  \approx 50.
\end{equation}
The simplified symmetric model yields $A \approx 50$.
Including the asymmetry and pairing corrections shifts the
peak to $A \approx 56$--$62$, consistent with the observed
$^{56}$Fe / $^{62}$Ni region.
The isotope $^{62}$Ni has the highest $B/A \approx
8.795$ MeV; $^{56}$Fe has $B/A \approx 8.790$ MeV and is
the most cosmically abundant due to its proximity to the
doubly-magic $^{56}$Ni ($Z = N = 28$).
\end{proof}

\begin{corollary}[Endpoint of Stellar Fusion]
\label{cor:fusion}
Fusion is exothermic for $A < A_{\mathrm{peak}}$ and
endothermic for $A > A_{\mathrm{peak}}$.  Stellar
nucleosynthesis terminates at the iron peak; the cores of
massive stars accumulate iron-group elements until
gravitational collapse.
\end{corollary}

\begin{remark}[$\varphi$-Ladder Proximity]
The iron peak $A_{\mathrm{peak}} \approx 56$ is close to
$\varphi^8 \approx 46.98$, within 20\%.  Whether this
proximity is structurally significant---connecting the peak
location to the 8th power of the golden ratio---or coincidental
is an open question.  The factorization $56 = 7 \times 8$
(an octave multiple of the 8-tick epoch) is suggestive but
not yet derived from the ledger structure.
\end{remark}

%% ============================================================
\section{Comparison with Experiment}
\label{sec:experiment}
%% ============================================================

\begin{center}
\renewcommand{\arraystretch}{1.3}
\begin{tabular}{@{}lccc@{}}
\toprule
\textbf{Nucleus} & \textbf{$B/A$ (expt.)} &
\textbf{$B/A$ (BW)} & \textbf{Notes} \\
\midrule
$^4$He & 7.07 MeV & 6.82 MeV & Doubly-magic ($Z = N = 2$) \\
$^{12}$C & 7.68 MeV & 7.56 MeV & Triple-$\alpha$ product \\
$^{16}$O & 7.98 MeV & 7.89 MeV & Doubly-magic ($Z = N = 8$) \\
$^{56}$Fe & 8.79 MeV & 8.74 MeV & Near doubly-magic \\
$^{62}$Ni & 8.79 MeV & 8.75 MeV & Highest $B/A$ \\
$^{208}$Pb & 7.87 MeV & 7.84 MeV & Doubly-magic
($Z=82$, $N=126$) \\
$^{238}$U & 7.57 MeV & 7.55 MeV & Fissile \\
\bottomrule
\end{tabular}
\end{center}

The BW formula with the five empirical coefficients reproduces
$B/A$ to within $\sim 1\%$ for most nuclei.  The RS
contribution is to provide a \emph{physical interpretation} of
each term in the context of the $J$-cost functional and to
derive the qualitative features (iron peak, magic number
enhancement, $N \approx Z$ preference) from ledger topology.

%% ============================================================
\section{Magic Numbers and Shell Structure}
\label{sec:magic}
%% ============================================================

\begin{definition}[Nuclear Magic Numbers]
\label{def:magic}
The nuclear magic numbers are the values
$Z,N \in \{2, 8, 20, 28, 50, 82, 126\}$ at which nuclei
exhibit enhanced stability.
\end{definition}

\begin{proposition}[Magic Numbers from Ledger Closures]
\label{prop:magic}
In RS, magic numbers correspond to ledger-neutral shell
closures.  The first two ($2, 8$) follow from the 8-tick
structure; higher magic numbers reflect 3D spherical packing
with spin-orbit splitting.  A first-principles derivation of
the complete sequence from the RS Hamiltonian is an open
problem.
\end{proposition}

\begin{remark}
The nuclear magic sequence shares its first two terms ($2, 8$)
with the electronic noble-gas sequence.  This is not a
coincidence in RS: both arise from the 8-tick phase structure.
The divergence at higher numbers ($20, 28$ vs.\ $18, 36$)
reflects the different packing geometry (3D sphere vs.\ radial
filling).
\end{remark}

%% ============================================================
\section{Open Problems}
\label{sec:open}
%% ============================================================

\begin{enumerate}
  \item \textbf{First-principles coefficients}: Deriving $a_V$,
  $a_S$, $a_A$, and $\delta_0$ from the $\varphi$-ladder
  without empirical input.  This requires a complete RS
  treatment of QCD confinement and the residual nuclear force.

  \item \textbf{Shell model from ledger}: Deriving the full
  magic number sequence from the 3D ledger topology, including
  spin-orbit splitting.

  \item \textbf{Drip lines}: Predicting the proton and neutron
  drip lines (limits of nuclear stability) from the $J$-cost
  framework.

  \item \textbf{Superheavy elements}: Predicting the island
  of stability near $Z = 114$, $N = 184$ from $\varphi$-ladder
  shell closures.
\end{enumerate}

%% ============================================================
\section{Predictions and Falsifiers}
\label{sec:falsifiers}
%% ============================================================

\begin{enumerate}[label=(\roman*)]
  \item \textbf{Asymmetry from $J$-cost}: The $(A-2Z)^2/A$
  dependence of the asymmetry term is a direct consequence of
  $J(N/Z)$.  Any deviation from this quadratic form at
  leading order would falsify the $J$-cost interpretation.

  \item \textbf{Iron peak location}: $A_{\mathrm{peak}} \in
  [50, 65]$, determined by the volume-surface-Coulomb balance.

  \item \textbf{Pairing sign}: Even-even nuclei must be more
  stable than their odd-$A$ neighbors.  This follows from
  $J$-cost minimization at $x = 1$ for paired nucleons.

  \item \textbf{Magic number enhancement}: Nuclei with
  $Z$ or $N \in \{2, 8, 20, 28, 50, 82, 126\}$ must show
  enhanced $B/A$ relative to neighbors.

  \item \textbf{Normal ordering}: The coefficient hierarchy
  $a_A > a_S > a_V \gg a_C$ must hold.  Reversal would
  contradict the relative magnitudes of the recognition
  filling, boundary, and electromagnetic scales.
\end{enumerate}

%% ============================================================
\section{Conclusion}
\label{sec:conclusion}
%% ============================================================

The Bethe--Weizs\"acker formula's functional form emerges
naturally from the RS framework:
\begin{enumerate}
  \item The volume term ($\propto A$) from recognition filling.
  \item The surface term ($\propto A^{2/3}$) from boundary
  recognition cost.
  \item The Coulomb term ($\propto Z^2/A^{1/3}$) from the
  derived $\alpha$.
  \item The asymmetry term ($\propto (A-2Z)^2/A$) directly
  from $J(N/Z)$.
  \item The pairing term from $\varphi$-resonant nucleon
  states.
\end{enumerate}
The iron peak, magic numbers, and nuclear stability emerge
from ledger topology.  The asymmetry term is the most direct
RS contribution: the BW asymmetry formula is literally the
Taylor expansion of the $J$-cost evaluated at the
neutron-proton ratio.  Deriving the numerical BW coefficients
from first principles remains an important open problem,
requiring a complete RS treatment of QCD confinement and
nuclear structure.

%% ============================================================
\begin{thebibliography}{99}

\bibitem{RS_Axioms}
J.~Washburn, ``The Algebra of Reality: A Recognition Science Derivation
of Physical Law,'' \emph{Axioms} \textbf{15}(2), 90 (2026).

\bibitem{BW1935}
C.~F.~von Weizs\"acker, ``Zur Theorie der Kernmassen,''
Z.\ Phys.\ \textbf{96}, 431 (1935).

\bibitem{Bethe1936}
H.~A.~Bethe and R.~F.~Bacher, ``Nuclear Physics A.
Stationary States of Nuclei,'' Rev.\ Mod.\ Phys.\
\textbf{8}, 82 (1936).

\bibitem{AME2020}
M.~Wang \textit{et al.}, ``The AME 2020 Atomic Mass
Evaluation,'' Chin.\ Phys.\ C \textbf{45}, 030003 (2021).

\bibitem{Krane}
K.~S.~Krane, \textit{Introductory Nuclear Physics}
(Wiley, 1988).

\end{thebibliography}

\end{document}
